\chapter{2026年全国I卷}

\section{单选题}

\begin{problem}
样本数据 \(6\)，\(8\)，\(4\)，\(5\)，\(12\) 的中位数为（\quad）

\choices
    {\(5\)}
    {\(6\)}
    {\(8\)}
    {\(9\)}

\end{problem}

\begin{answer}
B
\end{answer}

\begin{solution}
将这组数据按从小到大的顺序排列，得 \(4\)，\(5\)，\(6\)，\(8\)，\(12\). 共 \(5\) 个数，中间那个数是 \(6\)，所以中位数为 \(6\). 故选：B.
\end{solution}
\begin{problem}
已知平面向量 \(\bs{a}\)，\(\bs{b}\) 不共线，且 \(2\bs{a}+y\bs{b}=x\bs{a}-3\bs{b}\)，则（\quad）

\choices
    {\(x=2\)，\(y=-3\)}
    {\(x=-2\)，\(y=3\)}
    {\(x=2\)，\(y=3\)}
    {\(x=-2\)，\(y=-3\)}

\end{problem}

\begin{answer}
A
\end{answer}

\begin{solution}
因为 \(\bs{a}\)，\(\bs{b}\) 不共线，所以它们可以作为平面向量的一组基底，表示是唯一的. 比较
\(2\bs{a}+y\bs{b}=x\bs{a}-3\bs{b}\)
两边 \(\bs{a}\)，\(\bs{b}\) 的系数，得
\(x=2\)，\(y=-3\).
故选：A.
\end{solution}
\begin{problem}
已知集合 \(A=\left\{ \sin \frac{7\pi}{6},\cos \frac{5\pi}{3},\tan \frac{5\pi}{4} \right\}\)，\(B=\left\{ -\frac{\sqrt{3}}{2},-\frac{1}{2},1 \right\}\)，则 \(A\cap B=\)（\quad）

\choices
    {\(\left\{ -\frac{\sqrt{3}}{2},-\frac{1}{2} \right\}\)}
    {\(\left\{ -\frac{\sqrt{3}}{2},1 \right\}\)}
    {\(\left\{ -\frac{1}{2},1 \right\}\)}
    {\(\left\{ -\frac{\sqrt{3}}{2},-\frac{1}{2},1 \right\}\)}

\end{problem}

\begin{answer}
C
\end{answer}

\begin{solution}
因为 \(\sin \frac{7\pi}{6}=-\frac{1}{2}\)，\(\cos \frac{5\pi}{3}=\frac{1}{2}\)，\(\tan \frac{5\pi}{4}=1\)，所以 \(A=\left\{ -\frac{1}{2},\frac{1}{2},1 \right\}\). 于是 \(A\cap B=\left\{ -\frac{1}{2},1 \right\}\). 故选：C.
\end{solution}
\begin{problem}
曲线 \(y=5x+8\ln x\) 在点 \((1,5)\) 处的切线方程为（\quad）

\choices
    {\(y=3x+2\)}
    {\(y=5x\)}
    {\(y=8x-3\)}
    {\(y=13x-8\)}

\end{problem}

\begin{answer}
D
\end{answer}

\begin{solution}
函数 \(y=5x+8\ln x\) 的导数为 \(y'=5+\frac{8}{x}\). 当 \(x=1\) 时，切线斜率为 \(y'(1)=5+8=13\). 所以过点 \((1,5)\) 的切线方程为 \(y-5=13\left( x-1 \right)\)，即 \(y=13x-8\). 故选：D.
\end{solution}
\begin{problem}
已知抛物线 \(C_1:y^2=2p_1x\)（\(p_1>0\)）和 \(C_2:x^2=2p_2y\)（\(p_2>0\)）均经过点 \((4,8)\)，则 \(C_1\) 的焦点与 \(C_2\) 的焦点之间的距离为（\quad）

\choices
    {\(12\)}
    {\(4\sqrt{5}\)}
    {\(6\)}
    {\(\dfrac{\sqrt{65}}{2}\)}

\end{problem}

\begin{answer}
D
\end{answer}

\begin{solution}
点 \((4,8)\) 在 \(C_1\) 上，所以 \(8^2=2p_1\cdot 4\)，解得 \(p_1=8\). 故 \(C_1\) 的焦点为 \(\left( 4,0 \right)\).

点 \((4,8)\) 在 \(C_2\) 上，所以 \(4^2=2p_2\cdot 8\)，解得 \(p_2=1\). 故 \(C_2\) 的焦点为 \(\left( 0,\frac{1}{2} \right)\).

于是两焦点间距离为 \(\sqrt{\left( 4-0 \right)^2+\left( 0-\frac{1}{2} \right)^2}=\sqrt{16+\frac{1}{4}}=\frac{\sqrt{65}}{2}\). 故选：D.
\end{solution}
\begin{problem}
已知函数 \(f(x)=\dfrac{x+2}{\e^x+a}\) 的最大值为 \(1\)，则 \(a=\)（\quad）

\choices
    {\(\dfrac{1}{2}\)}
    {\(1\)}
    {\(\dfrac{3}{2}\)}
    {\(2\)}

\end{problem}

\begin{answer}
B
\end{answer}

\begin{solution}
因为 \(f(x)\) 的最大值为 \(1\)，所以对任意 \(x\in\R\)，都有 \(\frac{x+2}{\e^x+a}\le 1\). 于是 \(x+2\le \e^x+a\)，即 \(a\ge x+2-\e^x\). 设 \(g(x)=x+2-\e^x\)，则 \(g'(x)=1-\e^x\). 当 \(x<0\) 时，\(g'(x)>0\)；当 \(x>0\) 时，\(g'(x)<0\). 所以 \(g(x)\) 在 \(x=0\) 处取得最大值，且 \(g(0)=0+2-1=1\). 故 \(a\ge 1\).

若 \(a>1\)，则对任意 \(x\)，都有 \(x+2<\e^x+a\)，从而 \(f(x)<1\)，这与最大值为 \(1\) 矛盾. 因此只能有 \(a=1\). 故选：B.
\end{solution}
\begin{problem}
一百零八塔位于宁夏回族自治区青铜峡市，以其独特的建筑格局和深邃的历史文化闻名遐迩. 该塔群共有 \(108\) 座塔，依山势自上而下排成 \(12\) 行，将第 \(i\) 行中塔的座数记为 \(a_i\)（\(i=1,2,\cdots,12\)），其中 \(a_1=1\)，\(a_2=a_3=3\)，\(a_4=a_5=5\)，且 \(a_6,a_7,\cdots,a_{12}\) 是一个首项为 \(7\)、公差为 \(2\) 的等差数列. 将 \(a_1,a_2,\cdots,a_{12}\) 分为 \(6\) 组，每组 \(2\) 个数，使得每组的 \(2\) 个数之和可构成一个项数为 \(6\) 且公差为 \(d\)（\(d>0\)）的等差数列，则 \(d=\)（\quad）

\choices
    {\(2\)}
    {\(4\)}
    {\(6\)}
    {\(8\)}

\end{problem}

\begin{answer}
B
\end{answer}

\begin{solution}
由题意可得 \(a_1,a_2,\cdots,a_{12}=1,3,3,5,5,7,9,11,13,15,17,19\). 把这 \(12\) 个数分成 \(6\) 组后，六个组和的总和仍为 \(1+3+3+5+5+7+9+11+13+15+17+19=108\). 因此这 \(6\) 个组和所成等差数列的平均数为 \(\frac{108}{6}=18\).

又因为每个 \(a_i\) 都是奇数，所以每组两个数之和都是偶数. 设这 \(6\) 个组和从小到大为 \(x\)，\(x+d\)，\(x+2d\)，\(x+3d\)，\(x+4d\)，\(x+5d\). 由平均数为 \(18\)，得 \(x+\frac{5d}{2}=18\).

依次代入选项检验：

当 \(d=2\) 时，\(x=13\)，是奇数，不可能；

当 \(d=6\) 时，\(x=3\)，是奇数，不可能；

当 \(d=8\) 时，\(x=-2\)，不可能作为两数组和；

所以只能有 \(d=4\).

再给出一种实际分组：\(\left( 1,7 \right)\)，\(\left( 3,9 \right)\)，\(\left( 3,13 \right)\)，\(\left( 5,15 \right)\)，\(\left( 5,19 \right)\)，\(\left( 11,17 \right)\). 对应组和为 \(8\)，\(12\)，\(16\)，\(20\)，\(24\)，\(28\)，它们确实构成公差为 \(4\) 的等差数列. 故选：B.
\end{solution}
\begin{problem}
设 \(U=\{(x_1,x_2,x_3)\mid x_i\in\{-2,-1,1,2\},i=1,2,3\}\) 为空间中 \(64\) 个点构成的集合. 点 \(P(1,1,1)\)，记样本空间 \(\Omega=\complement_U \{P\}\). 从 \(\Omega\) 中随机取一个点，定义随机变量 \(X\) 如下：对 \(\Omega\) 中的每个点 \(A(x_1,x_2,x_3)\)，令 \(X(A)=x_1+x_2+x_3\)，则 \(X\) 的数学期望为（\quad）

\choices
    {\(-\dfrac{1}{21}\)}
    {\(-\dfrac{1}{63}\)}
    {\(0\)}
    {\(\dfrac{1}{7}\)}

\end{problem}

\begin{answer}
A
\end{answer}

\begin{solution}
在集合 \(U\) 中，\(x_1,x_2,x_3\) 都从 \(\{-2,-1,1,2\}\) 中取值. 由于这四个数的和为 \(-2+(-1)+1+2=0\)，所以在全部 \(64\) 个点中，\(x_1+x_2+x_3\) 的总和为 \(0\).

现在从 \(U\) 中删去点 \(P(1,1,1)\). 该点对应的函数值为 \(X(P)=1+1+1=3\). 因此在样本空间 \(\Omega\) 的 \(63\) 个点上，随机变量 \(X\) 的总和为 \(0-3=-3\). 于是数学期望为 \(E(X)=\frac{-3}{63}=-\frac{1}{21}\). 故选：A.
\end{solution}
\section{多选题}

\begin{problem}
设 \(z=3+2\i\)，则（\quad）

\choices
    {\(\overline{z}=3-2\i\)}
    {\(\lvert z \rvert=5\)}
    {\(z^2=5+12\i\)}
    {\(\dfrac{z+3}{z-\i}\in\R\)}

\end{problem}

\begin{answer}
ACD
\end{answer}

\begin{solution}
对于 A，\(\overline{z}=3-2\i\)，所以 A 正确.

对于 B，\(\lvert z \rvert=\sqrt{3^2+2^2}=\sqrt{13}\ne 5\)，所以 B 错误.

对于 C，\(z^2=\left( 3+2\i \right)^2=9+12\i+4\i^2=5+12\i\)，所以 C 正确.

对于 D，
\[
\frac{z+3}{z-\i}=\frac{6+2\i}{3+\i}=\frac{\left( 6+2\i \right)\left( 3-\i \right)}{\left( 3+\i \right)\left( 3-\i \right)}=\frac{20}{10}=2\in\R
\]
所以 D 正确.

故选：ACD.
\end{solution}
\begin{problem}
在空间中，\(A\)，\(B\) 为两个定点，动点 \(C\) 到直线 \(AB\) 的距离为 \(2\)，动点 \(D\) 到直线 \(AB\) 的距离为 \(1\)，若二面角 \(C-AB-D\) 为 \(60^\circ\)，则（\quad）

\choices
    {\(\angle CAD\ge 60^\circ\)}
    {\(CD\ge \sqrt{3}\)}
    {当 \(AB\perp CD\) 时，\(CD\perp \text{平面} ABD\)}
    {当 \(AB\perp \text{平面} ACD\) 时，\(AC\perp AD\)}

\end{problem}

\begin{answer}
BC
\end{answer}

\begin{solution}
取直线 \(AB\) 为 \(z\) 轴，设 \(A(0,0,0)\)，\(C(2,0,m)\)，\(D\left( \frac{1}{2},\frac{\sqrt{3}}{2},n \right)\). 这样点 \(C\) 到直线 \(AB\) 的距离为 \(2\)，点 \(D\) 到直线 \(AB\) 的距离为 \(1\)，且二面角 \(C-AB-D\) 为 \(60^\circ\).

对于 B，\(\overrightarrow{CD}=\left( -\frac{3}{2},\frac{\sqrt{3}}{2},n-m \right)\)，所以
\[
CD^2=\left( -\frac{3}{2} \right)^2+\left( \frac{\sqrt{3}}{2} \right)^2+\left( n-m \right)^2=3+\left( n-m \right)^2\ge 3
\]
故 \(CD\ge \sqrt{3}\)，所以 B 正确.

对于 C，若 \(AB\perp CD\)，则 \(\overrightarrow{CD}\) 的 \(z\) 分量为 \(0\)，即 \(n=m\). 此时 \(\overrightarrow{CD}=\left( -\frac{3}{2},\frac{\sqrt{3}}{2},0 \right)\)，\(\overrightarrow{AD}=\left( \frac{1}{2},\frac{\sqrt{3}}{2},m \right)\). 于是 \(\overrightarrow{CD}\cdot \overrightarrow{AD}=-\frac{3}{4}+\frac{3}{4}+0=0\)，所以 \(CD\perp AD\). 又已知 \(CD\perp AB\)，而 \(AB\) 与 \(AD\) 相交于点 \(A\)，故 \(CD\perp \text{平面} ABD\).
所以 C 正确.

对于 A，取 \(m=n=1\)，则 \(\overrightarrow{AC}=(2,0,1)\)，\(\overrightarrow{AD}=\left( \frac{1}{2},\frac{\sqrt{3}}{2},1 \right)\). 于是
\[
\cos \angle CAD=\frac{\overrightarrow{AC}\cdot \overrightarrow{AD}}{\lvert AC\rvert \lvert AD\rvert}=\frac{2}{\sqrt{5}\sqrt{2}}=\frac{2}{\sqrt{10}}>\frac{1}{2}
\]
故 \(\angle CAD<60^\circ\). 所以 A 错误.

对于 D，若 \(AB\perp \text{平面} ACD\)，则 \(m=n=0\). 此时 \(\overrightarrow{AC}=(2,0,0)\)，\(\overrightarrow{AD}=\left( \frac{1}{2},\frac{\sqrt{3}}{2},0 \right)\)，从而 \(\overrightarrow{AC}\cdot \overrightarrow{AD}=1\ne 0\).
所以 \(AC\perp AD\) 不成立，D 错误.

故选：BC.
\end{solution}
\begin{problem}
已知圆 \(C_1:(x+1)^2+y^2=1\)，圆 \(C_2:(x-1)^2+y^2=1\)，圆 \(C_3:x^2+\left( y-\sqrt{3} \right)^2=1\)，直线 \(l:y=kx+b\) 与 \(C_1\)，\(C_2\)，\(C_3\) 均有两个交点，记 \(l\) 被 \(C_1\)，\(C_2\)，\(C_3\) 截得弦长分别为 \(s_1\)，\(s_2\)，\(s_3\)，则（\quad）

\choices
    {\(k\) 可以取任意实数}
    {满足 \(s_1=s_2=s_3\) 的直线 \(l\) 共有 \(3\) 条}
    {满足 \(s_1+s_2+s_3=3\) 的直线 \(l\) 多于 \(3\) 条}
    {当 \(b=0\) 时，\(s_1+s_2+s_3\) 的最大值为 \(\dfrac{2\sqrt{21}}{3}\)}

\end{problem}

\begin{answer}
BCD
\end{answer}

\begin{solution}
设三个圆的圆心分别为 \(A(-1,0)\)，\(B(1,0)\)，\(C(0,\sqrt{3})\). 三圆半径都为 \(1\). 直线 \(l:y=kx+b\) 到三个圆心的距离分别为 \(d_1=\frac{\left| b-k \right|}{\sqrt{k^2+1}}\)，\(d_2=\frac{\left| b+k \right|}{\sqrt{k^2+1}}\)，\(d_3=\frac{\left| b-\sqrt{3} \right|}{\sqrt{k^2+1}}\). 于是 \(s_i=2\sqrt{1-d_i^2}\)（\(i=1,2,3\)）.

先看 A. 取 \(k=\frac{1}{\sqrt{3}}\). 若直线与三个圆都有两个交点，则必须同时满足 \(d_1<1\)，\(d_2<1\)，\(d_3<1\). 这等价于 \(\left| b-\frac{1}{\sqrt{3}} \right|<\frac{2}{\sqrt{3}}\)，\(\left| b+\frac{1}{\sqrt{3}} \right|<\frac{2}{\sqrt{3}}\)，\(\left| b-\sqrt{3} \right|<\frac{2}{\sqrt{3}}\). 化为区间分别是 \(\left( -\frac{1}{\sqrt{3}},\sqrt{3} \right)\)，\(\left( -\sqrt{3},\frac{1}{\sqrt{3}} \right)\)，\(\left( \frac{1}{\sqrt{3}},\frac{5}{\sqrt{3}} \right)\).
三者没有公共部分，所以当 \(k=\frac{1}{\sqrt{3}}\) 时不存在符合条件的直线. 因此 A 错误.

再看 B. 条件 \(s_1=s_2=s_3\) 等价于 \(d_1=d_2=d_3\).

由 \(d_1=d_2\) 可知，直线 \(l\) 或过原点，或平行于 \(x\) 轴.

若 \(l\) 平行于 \(x\) 轴，则 \(l:y=b\). 由 \(\left| b \right|=\left| b-\sqrt{3} \right|\) 得 \(b=\frac{\sqrt{3}}{2}\)，得到直线 \(y=\frac{\sqrt{3}}{2}\).

若 \(l\) 过原点，则 \(b=0\). 此时 \(d_1=d_3\) 化为 \(\frac{\left| k \right|}{\sqrt{k^2+1}}=\frac{\sqrt{3}}{\sqrt{k^2+1}}\)，故 \(\left| k \right|=\sqrt{3}\)，得到两条直线 \(y=\sqrt{3}x\)，\(y=-\sqrt{3}x\).
所以共有 \(3\) 条，B 正确.

再看 C. 取 \(b=0\). 此时为了与 \(C_3\) 有两个交点，需 \(k^2>2\). 并且 \(s_1=s_2=\frac{2}{\sqrt{k^2+1}}\)，\(s_3=\frac{2\sqrt{k^2-2}}{\sqrt{k^2+1}}\). 若 \(s_1+s_2+s_3=3\)，则 \(\frac{2\left( 2+\sqrt{k^2-2} \right)}{\sqrt{k^2+1}}=3\). 设 \(t=\sqrt{k^2-2}\)，则 \(t\ge 0\)，且 \(\sqrt{k^2+1}=\sqrt{t^2+3}\). 上式化为 \(2\left( 2+t \right)=3\sqrt{t^2+3}\).
两边平方得 \(5t^2-16t+11=0\)，所以 \(t=1\) 或 \(t=\frac{11}{5}\). 于是 \(k=\pm \sqrt{3}\) 或 \(k=\pm \frac{3\sqrt{19}}{5}\).
至少有 \(4\) 条直线满足条件，因此 C 正确.

最后看 D. 当 \(b=0\) 时，\(S=s_1+s_2+s_3=\frac{2\left( 2+\sqrt{k^2-2} \right)}{\sqrt{k^2+1}}\)（\(k^2>2\)）. 设 \(t=\sqrt{k^2-2}\)，则 \(S=\frac{2\left( 2+t \right)}{\sqrt{t^2+3}}\)，于是 \(S'=\frac{2\left( 3-2t \right)}{\left( t^2+3 \right)^{3/2}}\).
所以当 \(t=\frac{3}{2}\) 时，\(S\) 取得最大值. 此时
\[
S_{\max}=\frac{2\left( 2+\frac{3}{2} \right)}{\sqrt{\left( \frac{3}{2} \right)^2+3}}=\frac{7}{\sqrt{\frac{21}{4}}}=\frac{2\sqrt{21}}{3}
\]
故 D 正确.

故选：BCD.
\end{solution}
\section{填空题}

\begin{problem}
双曲线 \(5x^2-6y^2=1\) 的离心率为\fillinblank{}．

\end{problem}

\begin{answer}
\(\sqrt{\dfrac{11}{6}}\)
\end{answer}

\begin{solution}
将双曲线化为标准方程，得 \(\frac{x^2}{\frac{1}{5}}-\frac{y^2}{\frac{1}{6}}=1\). 所以 \(a^2=\frac{1}{5}\)，\(b^2=\frac{1}{6}\). 双曲线的离心率满足 \(e=\frac{c}{a}\)，\(c^2=a^2+b^2\). 因此
\[
e=\sqrt{1+\frac{b^2}{a^2}}=\sqrt{1+\frac{\frac{1}{6}}{\frac{1}{5}}}=\sqrt{1+\frac{5}{6}}=\sqrt{\frac{11}{6}}
\]
故填：\(\sqrt{\frac{11}{6}}\).
\end{solution}
\begin{problem}
已知 \(f(x)=2\sin(ax+\theta)\)（\(a\in\Z\)，\(0\le \theta<2\pi\)）是偶函数，\(f(x)\) 在区间 \(\left( 0,\frac{\pi}{2} \right)\) 单调递增，则 \(\theta=\fillinblank{}\)，\(f\left( \frac{2\pi}{3} \right)=\fillinblank{}\)．

\end{problem}

\begin{answer}
\(\dfrac{3\pi}{2}\)，\(1\)
\end{answer}

\begin{solution}
因为 \(f(x)\) 是偶函数，所以 \(f(-x)=f(x)\)，即 \(2\sin\left( -ax+\theta \right)=2\sin\left( ax+\theta \right)\). 将 \(f(x)\) 展开为 \(f(x)=2\sin\left( ax \right)\cos\theta+2\cos\left( ax \right)\sin\theta\). 因为偶函数中奇函数部分系数应为 \(0\)，所以 \(\cos\theta=0\). 结合 \(0\le \theta<2\pi\)，得 \(\theta=\frac{\pi}{2}\) 或 \(\theta=\frac{3\pi}{2}\).

若 \(\theta=\frac{\pi}{2}\)，则 \(f(x)=2\cos\left( ax \right)\). 若 \(a=0\)，则 \(f(x)\) 为常数函数，不符合在区间 \(\left( 0,\frac{\pi}{2} \right)\) 单调递增. 若 \(a\ne 0\)，其导数为 \(f'(x)=-2a\sin\left( ax \right)\).
当 \(0<|a|\le 2\) 时，\(f'(x)<0\)（\(0<x<\frac{\pi}{2}\)）；当 \(|a|\ge 3\) 时，\(\sin\left( ax \right)\) 在区间 \(\left( 0,\frac{\pi}{2} \right)\) 内变号，函数不可能在该区间单调递增. 所以这种情况不合题意.

故只能有 \(\theta=\frac{3\pi}{2}\). 此时 \(f(x)=-2\cos\left( ax \right)\). 若 \(a=0\)，则 \(f(x)\) 为常数函数，不符合题意. 若 \(a\ne 0\)，则 \(f'(x)=2a\sin\left( ax \right)\).
当 \(|a|=1\) 或 \(|a|=2\) 时，\(f'(x)>0\)（\(0<x<\frac{\pi}{2}\)），函数在区间 \(\left( 0,\frac{\pi}{2} \right)\) 单调递增；当 \(|a|\ge 3\) 时，\(\sin\left( ax \right)\) 在区间 \(\left( 0,\frac{\pi}{2} \right)\) 内变号，函数不可能在该区间单调递增.

因此满足题意的整数 \(a\) 为 \(a=\pm1\) 或 \(a=\pm2\). 这些情况下都有 \(f\left( \frac{2\pi}{3} \right)=1\). 例如当 \(a=1\) 时，\(f\left( \frac{2\pi}{3} \right)=-2\cos \frac{2\pi}{3}=1\)；当 \(a=2\) 时，\(f\left( \frac{2\pi}{3} \right)=-2\cos \frac{4\pi}{3}=1\)；当 \(a=-1\) 时，\(f\left( \frac{2\pi}{3} \right)=-2\cos \left( -\frac{2\pi}{3} \right)=1\)；当 \(a=-2\) 时，\(f\left( \frac{2\pi}{3} \right)=-2\cos \left( -\frac{4\pi}{3} \right)=1\).
故填：\(\frac{3\pi}{2}\)，\(1\).
\end{solution}
\begin{problem}
设实数 \(q\) 满足：存在数列 \(\{a_n\}\)，使得对于任意 \(n\in\N^*\)，均有 \(a_1+a_2+\cdots+a_{3n}=n^2+n\)，且 \(\{a_n\}\) 中有某连续 \(9\) 项 \(a_k,a_{k+1},\cdots,a_{k+8}\) 是公比为 \(q\) 的等比数列，则 \(q\) 的最大值为\fillinblank{}．

\end{problem}

\begin{answer}
\(\sqrt[3]{\dfrac{3}{2}}\)
\end{answer}

\begin{solution}
设 \(S_n=a_1+a_2+\cdots+a_n\). 由题意，\(S_{3n}=n^2+n\). 于是对任意 \(n\in\N^*\)，有 \(a_{3n-2}+a_{3n-1}+a_{3n}=S_{3n}-S_{3n-3}=n^2+n-\left( n-1 \right)n=2n\).
也就是说，每连续 \(3\) 项 \(a_{3n-2}\)，\(a_{3n-1}\)，\(a_{3n}\) 的和等于 \(2n\).

设某连续 \(9\) 项 \(a_t,a_{t+1},\cdots,a_{t+8}\) 是公比为 \(q\) 的等比数列，记 \(a_t=b\)，则这 \(9\) 项依次为 \(b\)，\(bq\)，\(bq^2\)，\(\cdots\)，\(bq^8\).

分三种情况讨论 \(t\) 对 \(3\) 的余数.

第一种，\(t=3m-2\). 这 \(9\) 项恰好分成三个完整的三项组，所以 \(b\left( 1+q+q^2 \right)=2m\)，\(bq^3\left( 1+q+q^2 \right)=2m+2\)，\(bq^6\left( 1+q+q^2 \right)=2m+4\). 于是 \(q^3=\frac{m+1}{m}\)，\(q^6=\frac{m+2}{m}\). 但这要求 \(\left( \frac{m+1}{m} \right)^2=\frac{m+2}{m}\).
化简得 \(1=0\)，矛盾. 所以这种情况不可能.

第二种，\(t=3m-1\). 这时其中完整的两个三项组分别是 \(a_{3m+1},a_{3m+2},a_{3m+3}\) 和 \(a_{3m+4},a_{3m+5},a_{3m+6}\). 因此 \(bq^2\left( 1+q+q^2 \right)=2m+2\)，\(bq^5\left( 1+q+q^2 \right)=2m+4\). 两式相除得 \(q^3=\frac{m+2}{m+1}\).

第三种，\(t=3m\). 这时其中完整的两个三项组分别是 \(a_{3m+1},a_{3m+2},a_{3m+3}\) 和 \(a_{3m+4},a_{3m+5},a_{3m+6}\). 对应到等比数列中的项为 \(bq,bq^2,bq^3\) 和 \(bq^4,bq^5,bq^6\). 因此这两个三项组满足 \(bq\left( 1+q+q^2 \right)=2m+2\)，\(bq^4\left( 1+q+q^2 \right)=2m+4\). 所以仍有 \(q^3=\frac{m+2}{m+1}\).

由第二，三种情况可知 \(q^3=\frac{m+2}{m+1}=1+\frac{1}{m+1}\)，故 \(q\) 随 \(m\) 的增大而减小，所以 \(q\le \sqrt[3]{\frac{3}{2}}\).

下面说明上界可以取得. 取 \(m=1\)，则 \(q=\sqrt[3]{\frac{3}{2}}\). 令 \(a_2,a_3,\cdots,a_{10}\) 成公比为 \(q\) 的等比数列，并令 \(a_2=b\)，\(a_3=bq\)，\(\cdots\)，\(a_{10}=bq^8\)，其中 \(b\) 取满足 \(bq^2\left( 1+q+q^2 \right)=4\).
这样便有 \(a_4+a_5+a_6=4\)，\(a_7+a_8+a_9=6\). 再取 \(a_1=2-a_2-a_3\)，\(a_{11}=0\)，\(a_{12}=8-a_{10}\)，
其余各组按“三项和等于 \(2n\)”任意补足即可，便能构造出符合题意的数列.

所以 \(q\) 的最大值为 \(\sqrt[3]{\frac{3}{2}}\). 故填：\(\sqrt[3]{\frac{3}{2}}\).
\end{solution}
\section{解答题}

\begin{problem}
如图，在直三棱柱 \(ABC-A_1B_1C_1\) 中，\(\angle ACB=90^\circ\)，\(AC=BC\)，\(D\)，\(E\) 分别为 \(AB\)，\(AC_1\) 的中点.

\begin{enumerate}
\item 证明：\(DE // \text{平面} BCC_1B_1\)；
\item 设 \(CC_1=2\)，直线 \(DE\) 与平面 \(ACC_1A_1\) 所成的角为 \(45^\circ\)，求直线 \(DE\) 到平面 \(BCC_1B_1\) 的距离．
\end{enumerate}

\bitmapfigure[max width=5.6cm,max totalheight=4.8cm]{img/2026/national_paper_1/q15_fig1.png}

\end{problem}

\begin{solution}
设 \(AC=BC=a\). 建立空间直角坐标系，使 \(C(0,0,0)\)，\(A(a,0,0)\)，\(B(0,a,0)\)，\(C_1(0,0,2)\). 因为是直三棱柱，所以 \(A_1(a,0,2)\)，\(B_1(0,a,2)\). 又因为 \(D\) 为 \(AB\) 的中点，\(E\) 为 \(AC_1\) 的中点，所以 \(D\left( \frac{a}{2},\frac{a}{2},0 \right)\)，\(E\left( \frac{a}{2},0,1 \right)\).

（1）平面 \(BCC_1B_1\) 由四点 \(B(0,a,0)\)，\(C(0,0,0)\)，\(C_1(0,0,2)\)，\(B_1(0,a,2)\) 确定，其方程为 \(x=0\). 而 \(\overrightarrow{DE}=\left( 0,-\frac{a}{2},1 \right)\). 向量 \(\overrightarrow{DE}\) 的 \(x\) 分量为 \(0\)，所以直线 \(DE\) 的方向向量与平面 \(x=0\) 平行，从而 \(DE // \text{平面} BCC_1B_1\).

（2）平面 \(ACC_1A_1\) 由 \(A(a,0,0)\)，\(C(0,0,0)\)，\(C_1(0,0,2)\)，\(A_1(a,0,2)\) 确定，其方程为 \(y=0\). 平面 \(y=0\) 的一个法向量为 \(\bs{n}=(0,1,0)\).

设直线 \(DE\) 与平面 \(ACC_1A_1\) 所成的角为 \(\theta\). 由题意 \(\theta=45^\circ\). 根据线面角公式，\(\sin\theta=\frac{\left| \overrightarrow{DE}\cdot \bs{n} \right|}{\lvert \overrightarrow{DE}\rvert}=\frac{\frac{a}{2}}{\sqrt{\frac{a^2}{4}+1}}\).
代入 \(\theta=45^\circ\)，得 \(\frac{\sqrt{2}}{2}=\frac{\frac{a}{2}}{\sqrt{\frac{a^2}{4}+1}}\). 两边平方，得 \(\frac{1}{2}=\frac{\frac{a^2}{4}}{\frac{a^2}{4}+1}\). 解得 \(a=2\).

由（1）知 \(DE // \text{平面} BCC_1B_1\)，所以直线 \(DE\) 到平面 \(BCC_1B_1\) 的距离等于点 \(D\) 到平面 \(x=0\) 的距离，即 \(\frac{a}{2}=1\). 所以所求距离为 \(1\).
\end{solution}
\begin{problem}
已知在 \(\triangle ABC\) 中，\(AB=3\)，\(BC=2\sqrt{3}\)，\(\cos B=\dfrac{\sqrt{3}}{3}\).

\begin{enumerate}
\item 求 \(\cos A\)；
\item 设 \(D\)，\(E\) 两点满足：\(D\) 在 \(BA\) 的延长线上，\(DE // BC\)，\(AE\perp AC\). 若 \(DE=\sqrt{6}\)，求 \(CE\)．
\end{enumerate}

\end{problem}

\begin{solution}
（1）由余弦定理，\(AC^2=AB^2+BC^2-2AB\cdot BC\cdot \cos B\). 代入已知，得 \(AC^2=3^2+\left( 2\sqrt{3} \right)^2-2\cdot 3\cdot 2\sqrt{3}\cdot \frac{\sqrt{3}}{3}=9+12-12=9\)，所以 \(AC=3\). 再由余弦定理，\(\cos A=\frac{AB^2+AC^2-BC^2}{2AB\cdot AC}=\frac{9+9-12}{2\cdot 3\cdot 3}=\frac{1}{3}\).

（2）取 \(A(0,0)\)，\(C(3,0)\). 设 \(B(x,y)\). 由 \(AB=3\)，\(BC=2\sqrt{3}\)，得 \(x^2+y^2=9\)，\(\left( x-3 \right)^2+y^2=12\). 两式相减，得 \(-6x+9=3\)，所以 \(x=1\). 再代回 \(y^2=9-1=8\). 取 \(B\left( 1,2\sqrt{2} \right)\).

因为 \(D\) 在 \(BA\) 的延长线上，所以可设 \(D=\lambda \left( -1,-2\sqrt{2} \right)\)（\(\lambda>0\)）. 又因为 \(AE\perp AC\)，而 \(AC\) 在 \(x\) 轴上，所以 \(E\) 在过 \(A\) 的 \(y\) 轴上，设 \(E(0,t)\).

由 \(DE // BC\)，得 \(\overrightarrow{DE}\ //\ \overrightarrow{BC}\). 其中 \(\overrightarrow{BC}=\left( 2,-2\sqrt{2} \right)\)，\(\overrightarrow{DE}=\left( \lambda,t+2\sqrt{2}\lambda \right)\). 故存在实数 \(\mu\)，使 \(\left( \lambda,t+2\sqrt{2}\lambda \right)=\mu \left( 2,-2\sqrt{2} \right)\). 比较 \(x\) 坐标，得 \(\lambda=2\mu\). 再比较 \(y\) 坐标，得 \(t+2\sqrt{2}\lambda=-2\sqrt{2}\mu\)，所以 \(t=-3\sqrt{2}\lambda\).

因为 \(\lvert BC\rvert =2\sqrt{3}\)，且 \(\mu=\frac{\lambda}{2}\)，所以 \(DE=\frac{\lambda}{2}\cdot 2\sqrt{3}=\lambda \sqrt{3}\). 由 \(DE=\sqrt{6}\)，得 \(\lambda \sqrt{3}=\sqrt{6}\)，所以 \(\lambda=\sqrt{2}\). 于是 \(t=-6\)，即 \(E(0,-6)\).

\[
CE=\sqrt{\left( 3-0 \right)^2+\left( 0+6 \right)^2}=\sqrt{45}=3\sqrt{5}
\]
所以 \(CE=3\sqrt{5}\).
\end{solution}
\begin{problem}
设整数 \(N\ge 2\). 某同学用一个球进行投篮练习，至多投篮 \(N\) 次，当且仅当投中 \(1\) 次时或 \(N\) 次均未投中时，停止练习. 设该同学每次投中的概率为 \(p\)（\(0<p<1\)），各次投中与否相互独立. 记 \(X\) 为停止练习时该同学的投篮次数.

\begin{enumerate}
\item 当 \(N=4\)，\(p=\dfrac{1}{3}\) 时，求 \(X\) 的分布列；
\item 设 \(k\)，\(m\) 均为自然数.
\begin{enumerate}
\item 当 \(k\le N-1\) 时，求 \(P(X>k)\)；
\item 当 \(k+m\le N-1\) 时，证明：\(P(X>k+m\mid X>k)=P(X>m)\)．
\end{enumerate}
\end{enumerate}

\end{problem}

\begin{solution}
（1）当 \(N=4\)，\(p=\frac{1}{3}\) 时，\(P\left( X=1 \right)=\frac{1}{3}\)，\(P\left( X=2 \right)=\frac{2}{3}\cdot \frac{1}{3}=\frac{2}{9}\)，\(P\left( X=3 \right)=\left( \frac{2}{3} \right)^2\cdot \frac{1}{3}=\frac{4}{27}\).

当 \(X=4\) 时，有两种情况：前三次都未中且第四次投中，或四次都未中. 所以
\[
P\left( X=4 \right)=\left( \frac{2}{3} \right)^3\cdot \frac{1}{3}+\left( \frac{2}{3} \right)^4=\frac{8}{81}+\frac{16}{81}=\frac{8}{27}
\]

故 \(X\) 的分布列为
\begin{center}
\begin{tabular}{c|cccc}
\(X\) & \(1\) & \(2\) & \(3\) & \(4\) \\
\hline
\(P\) & \(\frac{1}{3}\) & \(\frac{2}{9}\) & \(\frac{4}{27}\) & \(\frac{8}{27}\)
\end{tabular}
\end{center}

（2）（i）事件 \(X>k\) 表示前 \(k\) 次都没有投中. 因为各次投篮相互独立，且每次未中的概率为 \(1-p\)，所以 \(P\left( X>k \right)=\left( 1-p \right)^k\).

（2）由条件概率公式，\(P\left( X>k+m\mid X>k \right)=\frac{P\left( X>k+m,\ X>k \right)}{P\left( X>k \right)}\). 因为事件 \(X>k+m\) 蕴含 \(X>k\)，所以 \(P\left( X>k+m,\ X>k \right)=P\left( X>k+m \right)\). 再由（i）可得 \(P\left( X>k+m \right)=\left( 1-p \right)^{k+m}\)，\(P\left( X>k \right)=\left( 1-p \right)^k\). 故 \(P\left( X>k+m\mid X>k \right)=\frac{\left( 1-p \right)^{k+m}}{\left( 1-p \right)^k}=\left( 1-p \right)^m\). 而由（i）又有 \(P\left( X>m \right)=\left( 1-p \right)^m\). 所以 \(P\left( X>k+m\mid X>k \right)=P\left( X>m \right)\).
\end{solution}
\begin{problem}
已知椭圆 \(C:\dfrac{x^2}{a^2}+\dfrac{y^2}{b^2}=1\)（\(a>b>0\)）的左焦点为 \(F(-1,0)\)，离心率为 \(\dfrac{1}{2}\).

\begin{enumerate}
\item 求 \(C\) 的方程；
\item 设 \(O\) 为坐标原点，过 \(F\) 且斜率大于 \(0\) 的动直线 \(l\) 与 \(C\) 交于 \(P\)，\(Q\) 两点，其中 \(Q\) 在第三象限，直线 \(PO\) 与 \(C\) 的另一个交点为 \(R\).
\begin{enumerate}
\item 若 \(\triangle PQR\) 的面积是 \(\triangle PFO\) 的面积的 \(3\) 倍，求 \(l\) 的方程；
\item 求 \(\tan \angle PQR\) 的最小值．
\end{enumerate}
\end{enumerate}

\end{problem}

\begin{solution}
（1）由左焦点为 \(F(-1,0)\)，得 \(c=1\). 又因为离心率 \(e=\frac{c}{a}=\frac{1}{2}\)，所以 \(a=2\). 于是 \(b^2=a^2-c^2=4-1=3\). 故椭圆 \(C\) 的方程为 \(\frac{x^2}{4}+\frac{y^2}{3}=1\).

（2）设直线 \(l\) 的方程为 \(y=k\left( x+1 \right)\)（\(k>0\)）. 设 \(P(x_1,y_1)\)，\(Q(x_2,y_2)\)，其中 \(x_1>x_2\). 将直线方程代入椭圆方程，得 \(\frac{x^2}{4}+\frac{k^2\left( x+1 \right)^2}{3}=1\). 化简为 \(\left( 3+4k^2 \right)x^2+8k^2x+4k^2-12=0\).
所以 \(x_1+x_2=-\frac{8k^2}{3+4k^2}\)，\(x_1-x_2=\frac{12\sqrt{k^2+1}}{3+4k^2}\).

因为 \(O\) 是椭圆中心，且直线 \(PO\) 经过原点，所以它与椭圆的另一个交点是 \(P\) 关于原点的对称点，即 \(R(-x_1,-y_1)\).

（1）由 \(y_1=k\left( x_1+1 \right)\)，\(y_2=k\left( x_2+1 \right)\)，可得 \(S_{\triangle PQR}=\frac{1}{2}\left| \det\left( \overrightarrow{QP},\overrightarrow{QR} \right) \right|=k\left( x_1-x_2 \right)\).

又因为 \(F(-1,0)\)，\(O(0,0)\)，所以 \(S_{\triangle PFO}=\frac{1}{2}\left| y_1 \right|=\frac{1}{2}k\left( x_1+1 \right)\).

由题意 \(S_{\triangle PQR}=3S_{\triangle PFO}\)，得 \(k\left( x_1-x_2 \right)=\frac{3}{2}k\left( x_1+1 \right)\). 约去 \(k\)，得 \(2\left( x_1-x_2 \right)=3\left( x_1+1 \right)\). 代入 \(x_1=\frac{-4k^2+6\sqrt{k^2+1}}{3+4k^2}\) 和 \(x_1-x_2=\frac{12\sqrt{k^2+1}}{3+4k^2}\)，得到 \(\frac{24\sqrt{k^2+1}}{3+4k^2}=\frac{9+18\sqrt{k^2+1}}{3+4k^2}\).
所以 \(6\sqrt{k^2+1}=9\)，即 \(\sqrt{k^2+1}=\frac{3}{2}\). 故 \(k^2=\frac{5}{4}\)，\(k=\frac{\sqrt{5}}{2}\). 于是直线 \(l\) 的方程为 \(y=\frac{\sqrt{5}}{2}\left( x+1 \right)\).

（2）设 \(\overrightarrow{QP}=P-Q\)，\(\overrightarrow{QR}=R-Q\). 因为 \(P,Q\) 都在直线 \(l\) 上，所以 \(\overrightarrow{QP}\) 可取方向向量 \(\left( 1,k \right)\). 又因为 \(R=-P\)，所以 \(\overrightarrow{QR}=R-Q=-(P+Q)\). 而 \(P+Q=\left( x_1+x_2,\ y_1+y_2 \right)=\left( x_1+x_2,\ k\left( x_1+x_2+2 \right) \right)\).

于是 \(\tan \angle PQR=\frac{\left| \det\left( \left( 1,k \right),\overrightarrow{QR} \right) \right|}{\left( 1,k \right)\cdot \overrightarrow{QR}}\).
计算得 \(\left| \det\left( \left( 1,k \right),\overrightarrow{QR} \right) \right|=2k\)，并且 \(\left( 1,k \right)\cdot \overrightarrow{QR}=\frac{2k^2}{3+4k^2}\).
因此
\[
\tan \angle PQR=\frac{2k}{\frac{2k^2}{3+4k^2}}=\frac{3+4k^2}{k}=4k+\frac{3}{k}
\]

由基本不等式，\(4k+\frac{3}{k}\ge 2\sqrt{12}=4\sqrt{3}\)，当且仅当 \(4k=\frac{3}{k}\)，即 \(k=\frac{\sqrt{3}}{2}\) 时取等号.

所以 \(\tan \angle PQR\) 的最小值为 \(4\sqrt{3}\).
\end{solution}
\begin{problem}
已知函数 \(f(x)\) 的定义域为 \(\R\)，且当 \(x<0\) 时，\(f(x)=2^x\). 对任意 \(x_0\in\R\)，定义集合 \(D(x_0)=\{d\in\R \mid f(x_0+d)>f(x_0)\}\).

\begin{enumerate}
\item 若当 \(x\ge 0\) 时，\(f(x)=1-x\)，求 \(D(-1)\)；
\item 若 \(f(x)\) 是奇函数，\(f(x_1)\le f(x_2)\)，且 \(x_1x_2\ne 0\)，证明：\(D(x_2)\subseteq D(x_1)\)；
\item 设 \(f(x)\) 满足：\circled{1} 若 \(f(x_1)\le f(x_2)\)，则 \(D(x_2)\subseteq D(x_1)\)；\circled{2} 当 \(0<x<1\) 时，\(f(x)<f(0)\).
\begin{enumerate}
\item 证明：\(f(0)\ge 1\)；
\item 证明：\(f(x)\) 在区间 \((0,+\infty)\) 单调递增．
\end{enumerate}
\end{enumerate}

\end{problem}

\begin{solution}
（1）因为 \(-1<0\)，所以 \(f(-1)=2^{-1}=\frac{1}{2}\). 由定义，\(D(-1)=\left\{ d\in\R \mid f(d-1)>\frac{1}{2} \right\}\).

当 \(d<1\) 时，\(d-1<0\)，所以 \(f(d-1)=2^{d-1}\). 由 \(2^{d-1}>\frac{1}{2}=2^{-1}\) 得 \(d>0\).
故此时 \(0<d<1\).

当 \(d\ge 1\) 时，\(d-1\ge 0\)，所以 \(f(d-1)=1-\left( d-1 \right)=2-d\). 由 \(2-d>\frac{1}{2}\) 得 \(d<\frac{3}{2}\).
故此时 \(1\le d<\frac{3}{2}\).

综上，\(D(-1)=\left( 0,\frac{3}{2} \right)\).

（2）因为 \(f(x)\) 是奇函数，且当 \(x<0\) 时 \(f(x)=2^x\)，所以当 \(x>0\) 时 \(f(x)=-f(-x)=-2^{-x}\). 于是
\[
f(x)=
\begin{cases}
2^x, & x<0, \\
0, & x=0, \\
-2^{-x}, & x>0.
\end{cases}
\]

先求出各点对应的 \(D\) 集合.

当 \(x<0\) 时，若 \(d\le 0\)，则 \(x+d<x\)，从而 \(f(x+d)=2^{x+d}\le 2^x=f(x)\). 所以 \(d\notin D(x)\). 若 \(0<d<-x\)，则 \(x+d<0\)，且 \(2^{x+d}>2^x\). 故 \(d\in D(x)\). 若 \(d\ge -x\)，则 \(x+d\ge 0\)，此时 \(f(x+d)\le 0<2^x=f(x)\). 故 \(d\notin D(x)\). 所以 \(D(x)=\left( 0,-x \right)\)（\(x<0\)）.

当 \(x>0\) 时，若 \(d>0\)，则 \(x+d>0\)，且 \(f(x+d)=-2^{-(x+d)}>-2^{-x}=f(x)\). 故 \(d\in D(x)\).

若 \(d=-x\)，则 \(x+d=0\)，于是 \(f(0)=0>-2^{-x}=f(x)\). 故 \(-x\in D(x)\).

若 \(d<-x\)，则 \(x+d<0\)，从而 \(f(x+d)=2^{x+d}>0>-2^{-x}=f(x)\). 故 \(d\in D(x)\).

若 \(-x<d<0\)，则 \(x+d>0\)，且 \(f(x+d)=-2^{-(x+d)}<-2^{-x}=f(x)\). 故 \(d\notin D(x)\). 所以 \(D(x)=\left( -\infty,-x \right]\cup \left( 0,+\infty \right)\)（\(x>0\)）.

下面分情形证明.

若 \(x_1<x_2<0\)，则 \(f(x_1)=2^{x_1}\le 2^{x_2}=f(x_2)\)，且 \(D(x_2)=\left( 0,-x_2 \right)\subseteq \left( 0,-x_1 \right)=D(x_1)\).

若 \(0<x_1<x_2\)，则 \(f(x_1)=-2^{-x_1}\le -2^{-x_2}=f(x_2)\)，且 \(\left( -\infty,-x_2 \right]\cup \left( 0,+\infty \right)\subseteq \left( -\infty,-x_1 \right]\cup \left( 0,+\infty \right)\). 所以 \(D(x_2)\subseteq D(x_1)\).

若 \(x_1<0<x_2\)，则 \(f(x_1)=2^{x_1}>0>-2^{-x_2}=f(x_2)\)，
这与 \(f(x_1)\le f(x_2)\) 矛盾，所以这种情况不可能出现.

若 \(x_2<0<x_1\)，则 \(f(x_2)=2^{x_2}>0>-2^{-x_1}=f(x_1)\)，从而 \(f(x_1)\le f(x_2)\) 恒成立. 此时 \(D(x_2)=\left( 0,-x_2 \right)\subseteq \left( 0,+\infty \right)\subseteq D(x_1)\).

综上，总有 \(D(x_2)\subseteq D(x_1)\).

（3）（i）反证法. 若 \(f(0)<1\)，则可取 \(x_0\in(-1,0)\)，使 \(f(0)<f(x_0)=2^{x_0}<1\). 再取任意 \(d\in\left( 0,-x_0 \right)\). 因为 \(x_0+d\in(x_0,0)\subset(-1,0)\)，且 \(2^x\) 在 \((-\infty,0)\) 上单调递增，所以 \(f(x_0+d)>f(x_0)\)，即 \(d\in D(x_0)\).

又因为 \(f(0)<f(x_0)\)，由条件 \circled{1} 得 \(D(x_0)\subseteq D(0)\). 故 \(d\in D(0)\)，于是 \(f(d)>f(0)\). 但 \(0<d<1\)，由条件 \circled{2} 应有 \(f(d)<f(0)\)，
矛盾.

所以 \(f(0)\ge 1\).

（2）先证当 \(0<x<1\) 时，\(f(x)\le 0\).

若存在 \(x_0\in(0,1)\) 使 \(f(x_0)>0\)，则由条件 \circled{2} 有 \(f(x_0)<f(0)\)，所以 \(-x_0\in D(x_0)\). 又可取足够小的负数 \(x_1<0\)，使 \(f(x_1)=2^{x_1}<f(x_0)\). 由条件 \circled{1} 得 \(D(x_0)\subseteq D(x_1)\). 于是 \(-x_0\in D(x_1)\)，即 \(f(x_1-x_0)>f(x_1)\).
但 \(x_1-x_0<x_1<0\)，而 \(2^x\) 在 \((-\infty,0)\) 上单调递增，故 \(f(x_1-x_0)=2^{x_1-x_0}<2^{x_1}=f(x_1)\).
矛盾.

所以当 \(0<x<1\) 时，必有 \(f(x)\le 0\).

再证当 \(x\ge 1\) 时，也有 \(f(x)\le 0\).

若存在 \(x_0\ge 1\) 使 \(f(x_0)>0\)，取 \(x_1<0\) 使 \(f(x_1)<f(x_0)\). 令 \(d=x_0-x_1>0\)，则 \(x_1+d=x_0\)，所以 \(d\in D(x_1)\).

再令 \(x_2=x_1-\left( x_0-\frac{1}{2} \right)\). 则 \(x_2<x_1<0\)，从而 \(f(x_2)=2^{x_2}<2^{x_1}=f(x_1)<f(x_0)\). 由条件 \circled{1} 得 \(D(x_1)\subseteq D(x_2)\). 故 \(d\in D(x_2)\). 于是 \(f(x_2+d)>f(x_2)\). 但 \(x_2+d=x_1-\left( x_0-\frac{1}{2} \right)+x_0-x_1=\frac{1}{2}\)，所以 \(f\left( \frac{1}{2} \right)>f(x_2)>0\).
这与前面已证的 \(0<x<1\) 时 \(f(x)\le 0\) 矛盾.

故对一切 \(x>0\)，都有 \(f(x)\le 0\).

最后证明单调递增. 任取 \(x_0>0\) 和 \(d>0\). 取整数 \(n>d\). 因为 \(f(x_0)\le 0<2^{-n}=f(-n)\)，由条件 \circled{1} 得 \(D(-n)\subseteq D(x_0)\). 又因为 \(-n< -n+d<0\)，且 \(2^x\) 在 \((-\infty,0)\) 上单调递增，所以 \(f(-n+d)>f(-n)\)，即 \(d\in D(-n)\). 从而 \(d\in D(x_0)\). 所以 \(f(x_0+d)>f(x_0)\).

这说明对任意 \(x_0>0\) 和任意 \(d>0\)，都有 \(f(x_0+d)>f(x_0)\).
于是 \(f(x)\) 在区间 \((0,+\infty)\) 上单调递增.
\end{solution}
